% Generated draft: requires mathematical and source review.
\documentclass[11pt]{article}
\usepackage[margin=25mm]{geometry}
\usepackage{fontspec}
\setmainfont{Latin Modern Roman}
\usepackage{amsmath,amssymb,mathtools,unicode-math}
\setmathfont{Latin Modern Math}
\usepackage{enumitem,needspace,xurl,hyperref}
\setlength{\parindent}{0pt}
\setlength{\parskip}{5pt}
\setlength{\emergencystretch}{3em}
\newcommand{\N}{\mathbb{N}}
\newcommand{\Z}{\mathbb{Z}}
\newcommand{\Q}{\mathbb{Q}}
\newcommand{\R}{\mathbb{R}}
\newcommand{\C}{\mathbb{C}}
\newcommand{\F}{\mathbb{F}}
\newcommand{\A}{\mathbb{A}}
\newcommand{\PP}{\mathbb{P}}
\newcommand{\E}{\mathbb{E}}
\newcommand{\eps}{\varepsilon}
\newcommand{\dd}{\,\mathrm{d}}
\begin{document}
\section*{UnsolvedMath ID: 700393}
\textbf{Source problem code:} catalogue entry
\subsection{Definitions and mathematical statement}
All groups under consideration are finite. A formation is a class of finite groups closed under homomorphic images and finite subdirect products. A subdirect product of groups \(G_1,\ldots,G_n\) is a subgroup of \(G_1\times\cdots\times G_n\) whose projection onto each factor is surjective.

Let \(\mathfrak X\) be a formation of finite groups, and let \(\mathfrak F\subseteq\mathfrak X\) be a local subformation in the terminology of the cited reference. Write \(s_{\mathfrak F}\) for the object called the maximal inner local screen of \(\mathfrak F\) in the source, and put
\[
\Omega(\mathfrak F,\mathfrak X)
=\bigl\{s : s\text{ is a maximal homogeneous }\mathfrak X\text{-screen of }\mathfrak F\bigr\}.
\]
For arbitrary such \(\mathfrak F\) and \(\mathfrak X\), find a construction of the elements of \(\Omega(\mathfrak F,\mathfrak X)\) using \(s_{\mathfrak F}\), and determine what can be said about
\[
\bigl|\Omega(\mathfrak F,\mathfrak X)\bigr|.
\]
The supplied record refers to Shemetkov's book for the meanings of local subformation, homogeneous \(\mathfrak X\)-screen, maximality of such screens, and maximal inner local screen. Those definitions are not reproduced in the record. Consequently, the displayed set retains the source terminology rather than supplying an independently checkable definition. The notation \(s_{\mathfrak F}\) follows the source's reference to the maximal inner local screen; it does not independently establish existence or uniqueness. The requested scope of the cardinality conclusion is also unspecified.
\subsection{Short English statement}
Given a local subformation \(\mathfrak F\) of a formation \(\mathfrak X\) of finite groups, construct all maximal homogeneous \(\mathfrak X\)-screens of \(\mathfrak F\) using its maximal inner local screen, and investigate how many such screens there are. The specialized screen terminology requires the definitions in the cited book.
\subsection{Sources}
\begin{itemize}
\item Local UnsolvedMath registry, permanent ID \(700393\), exact source problem code \texttt{catalogue entry}, title \emph{Kourovka Notebook Problem 6.51}. The supplied record labels the problem open; this label is not an independently verified literature assessment.
\item \emph{The Kourovka Notebook (Unsolved Problems in Group Theory)}, Problem 6.51, proposed by L. A. Shemetkov, sixth issue (1978). The supplied background also identifies the twenty-first issue edition (2026), \texttt{arXiv:1401.0300v48}. These bibliographic details are supplied provenance and have not been independently checked.
\item L. A. Shemetkov, \emph{Formations of finite groups}, Moscow: Nauka, 1978, in Russian. The supplied statement explicitly refers to this book for definitions.
\item The supplied collection description credits E. I. Khukhro and V. D. Mazurov as editors of the Kourovka Notebook collection covering Issues 1--21 (1965--2026).
\end{itemize}
\end{document}
